% =========================================================================
% 논문 개요 (빠른 센서로 CES 빈칸 채우기) — 읽기 쉬운 국문 개요
% 빌드: xelatex overview_ko.tex   (두 번 돌려야 목차/참조가 맞는다)
% 원본: Notion 「논문 개요 (빠른 센서로 CES 빈칸 채우기, 2026-08-17)」
% 수치 출처: docs/paper/paper_numbers.json — main.tex / main_ko.tex와 동일
% =========================================================================
\documentclass[10pt]{article}

\usepackage[margin=0.9in]{geometry}
\usepackage{amsmath,amssymb}
\usepackage{booktabs}
\usepackage{tabularx}
\usepackage{array}
\usepackage{longtable}
\usepackage{xcolor}
\usepackage{xeCJK}
\xeCJKsetup{CJKspace=true}
\setCJKmainfont{Malgun Gothic}
\setCJKsansfont{Malgun Gothic}
\setCJKmonofont{Malgun Gothic}
\setmonofont{Consolas}
\usepackage[colorlinks=true,linkcolor=blue,urlcolor=blue]{hyperref}
\usepackage{titlesec}
\usepackage{enumitem}
\usepackage{framed}

\hypersetup{
  pdftitle={논문 개요 — 빠른 센서로 CES 빈칸 채우기},
  pdfauthor={이승상 (Seungsang Lee)},
}

\renewcommand{\contentsname}{목차}
\renewcommand{\tablename}{표}
\renewcommand{\today}{{\number\year}년 {\number\month}월 {\number\day}일}

% 긴 파일 경로(\texttt)와 붙임표 없는 토큰 때문에 줄이 넘치지 않도록 여유를 준다.
\emergencystretch=4em
\hbadness=10000
\sloppy

\setlist[itemize]{leftmargin=1.4em,itemsep=2pt,topsep=3pt}
\setlist[enumerate]{leftmargin=1.6em,itemsep=2pt,topsep=3pt}
\titlespacing*{\section}{0pt}{14pt}{6pt}
\titlespacing*{\subsection}{0pt}{10pt}{4pt}
\renewcommand{\arraystretch}{1.25}

% 자주 쓰는 기호
\newcommand{\Ti}{\texttt{T\_i}}
\newcommand{\Vrot}{\texttt{V\_rot}}
\newcommand{\seqv}{\texttt{seq\_v2}}
\newcommand{\bthree}{\texttt{b3k8}}
\newcommand{\yes}{\ensuremath{\bigcirc}}
\newcommand{\no}{\ensuremath{\times}}

% 회색 배경 상자 (Notion callout / 인용 대응)
\definecolor{shadecolor}{RGB}{240,244,250}
\newenvironment{note}{\begin{snugshade}\small}{\end{snugshade}}

% 표 안에서 쓰는 좌측 정렬 가변폭 열
\newcolumntype{L}{>{\raggedright\arraybackslash}X}
\newcolumntype{C}{>{\centering\arraybackslash}X}

\title{\textbf{논문 개요 — 빠른 센서로 CES 빈칸 채우기}\\[4pt]
       \large KSTAR 전하교환분광 결측 구간의 인과적 복원}
\author{이승상 (Seungsang Lee)}
\date{2026년 8월 17일}

\begin{document}
\maketitle

\begin{note}
\textbf{이 문서는 처음 보는 사람이 읽을 수 있도록 풀어 쓴 논문 개요다.}
수치는 전부 저장소의 \texttt{docs/paper/paper\_numbers.json}(실험 결과에서 자동 수집)에서
왔고 영문 \texttt{main.tex}·국문 \texttt{main\_ko.tex}와 일치한다.
근거 기록은 \texttt{THESIS\_RESULTS.md} §8v--§8ab, 실험 규칙은
\texttt{experiments/PREREGISTRATION\_W2.md}.
\end{note}

\tableofcontents
\vspace{6pt}

% =========================================================================
\section{먼저, 이 글에 나오는 말들}

논문 용어를 그대로 쓰면 읽기 어렵기 때문에, 이 문서에서 반복해서 나오는 말을 먼저
정리한다. \textbf{아래만 이해하면 나머지는 전부 읽힌다.}

\textbf{핵심}: CES(\Ti\,·\,\Vrot)는 \textbf{느리고 자주 빈다}. 빠른 진단 3종은
\textbf{빠르고 안 빈다}. 그런데 빠른 진단은 \Ti\,·\,\Vrot 를 \textbf{직접 재지 않는다}.

\subsection{비교 상대 (누구를 이겨야 하나)}

우리 모델과 겨루는 상대들이다. 위로 갈수록 약하고, 아래로 갈수록 강하다.

\begin{center}
\begin{tabularx}{\textwidth}{@{}l L c@{}}
\toprule
\textbf{상대} & \textbf{무엇을 보나} & \textbf{실시간} \\
\midrule
persistence & 마지막으로 잰 값을 그대로 답으로 씀 & \yes \\
국소 AR & 과거 CES 값 몇 개로 추세를 그림 & \yes \\
인과 GP & 과거 CES 이웃 16개로 통계적으로 매끄러운 곡선을 그림 & \yes \\
선형 보간 & 앞뒤 값을 직선으로 이음 & \no \\
PCHIP & 앞뒤 값을 곡선으로 이음 & \no \\
GP (오프라인) & 앞뒤 이웃 16+16개로 잡음까지 감안해 곡선을 그림. \textbf{가장 강한 상대} & \no \\
\bottomrule
\end{tabularx}
\end{center}

\textbf{GP(가우시안 프로세스)}는 융합 연구에서 표준으로 쓰는 곡선 맞춤 방법이다.
``이 데이터에는 잡음이 이만큼 있다''를 스스로 추정해서, 점을 정확히 통과하는 대신
\textbf{잡음을 평균 내며} 지나간다. 그래서 PCHIP보다 강하다.

\textbf{왜 일부러 불리하게 겨루나}: 보간·GP는 미래를 보는데 우리 모델은 과거만 본다.
이 불공평한 조건에서도 이긴다면, ``빠른 진단이 CES 시간축만으로는 절대 알 수 없는 정보를
실제로 담고 있다''가 증명된다.

\subsection{점수 읽는 법}

\begin{center}
\begin{tabularx}{\textwidth}{@{}l L@{}}
\toprule
\textbf{말} & \textbf{뜻} \\
\midrule
paired 비교 & 두 방법을 \textbf{완전히 같은 행, 같은 조건}에서 나란히 채점.
  조건 차이 때문에 생기는 착각을 없앰 \\
shot 군집 bootstrap & 신뢰구간 계산법. 데이터를 \textbf{방전 단위로} 무작위 재추출해
  1만 번 다시 채점하고, 그 분포의 95\,\% 구간이 0을 넘지 않으면 ``유의하다''고 판정.
  \textbf{행이 아니라 방전을 단위로 쓰는 이유}: 같은 방전 안의 이웃 행끼리는 서로 닮아서,
  행을 단위로 쓰면 실제보다 데이터가 많은 척하게 됨 \\
\bottomrule
\end{tabularx}
\end{center}

\subsection{데이터 처리 규칙 (두 가지가 헷갈리기 쉽다)}
\label{sec:cut}

\textbf{\textcircled{1} \Ti\ 이상값 --- ``컷''과 ``포함''}

CES가 빛의 폭을 맞출 때 가끔 \textbf{맞추기에 실패}한다. 그러면 \Ti 가 3{,}000\,eV를
훌쩍 넘는 말도 안 되는 값이 찍힌다(최대 14{,}984\,eV). 전체의 \textbf{0.53\,\%}(1{,}197행)다.

\begin{itemize}
\item 이걸 \textbf{빼고} 채점하는 데이터 = \textbf{``컷''}
\item 이걸 \textbf{넣고} 채점하는 데이터 = \textbf{``포함''}
\end{itemize}

빼면 ``어려운 문제를 골라 버렸다''는 비판을, 넣으면 ``물리값도 아닌 걸로 채점한다''는
비판을 받는다. 그래서 \textbf{둘 다 1차 결과로 보고}하고, \textbf{두 쪽에서 모두 성립할
때만 조건 없는 주장}으로 쓴다.

\textbf{\textcircled{2} 유지값 제거}

\Vrot\ 관측값의 \textbf{54\,\%}가 유지값(직전 값 복사)이다. 이건 측정이 아니므로
\textbf{학습에서도 채점에서도 전부 제거}한다. 제거하지 않으면 ``직전 값 그대로'' 방법이
공짜로 정답을 맞히게 되어 비교가 무의미해진다.

\subsection{모델 관련 말}

\begin{center}
\begin{tabularx}{\textwidth}{@{}l L@{}}
\toprule
\textbf{말} & \textbf{뜻} \\
\midrule
윈도(window) 방식 & 답을 채울 칸 \textbf{바로 앞 2줄만 잘라서} 보는 방식.
  우리의 \textbf{비교용 대조군} \\
전체격자 시퀀스 방식 & 그 방전의 \textbf{모든 칸을 시간 순서대로 전부} 읽는 방식.
  라벨(정답)이 없는 칸도 \textbf{맥락으로} 읽는다. 우리의 \textbf{주 모델} \\
\seqv & 그 주 모델의 이름. 35만 7천 개 숫자(파라미터)로 이루어짐 \\
분기(branch) & 모델 안의 갈래. \Ti 용 갈래와 \Vrot 용 갈래가 \textbf{따로} 있고,
  \Vrot\ 갈래는 빠른 진단을 \textbf{아예 못 보게 배선}되어 있다 \\
도달거리(reach) & 모델이 과거를 얼마나 멀리까지 볼 수 있는가.
  윈도 방식은 2줄, 시퀀스 방식은 방전 구간 전체 \\
절제(ablation) & 입력의 일부를 일부러 빼고 다시 학습·채점해서, 그 입력이 실제로
  기여하는지 확인하는 실험 \\
\bottomrule
\end{tabularx}
\end{center}

% =========================================================================
\section{이 연구가 답하려는 질문}

\subsection{상황}

핵융합 플라즈마의 가장자리 이온 온도와 회전 속도는 \textbf{성능을 좌우하는 핵심 물리량}이다.
고성능 운전 모드로 넘어갈지, ELM이라는 가장자리 폭발이 일어날지가 여기서 갈린다.

그런데 이걸 재는 CES는 \textbf{빛을 충분히 모아야 해서 느리고}, 신호가 약하면 그 시점을
통째로 놓친다. 우리 데이터에서 0.01초 간격 칸을 기준으로:

\begin{itemize}
\item \Ti: \textbf{8.2\,\%가 빔}
\item \Vrot: \textbf{23.9\,\%가 비고}, 거기에 더해 \textbf{41.1\,\%는 유지값}(직전 값 복사)
\end{itemize}

$\rightarrow$ \Vrot 는 전체 칸의 \textbf{35\,\%에만 진짜 새 정보}가 있다.

게다가 값이 빠지는 시점은 무작위가 아니다. \textbf{신호가 약해지는 순간, ELM이 터지는 순간,
모드가 바뀌는 순간} --- 즉 물리적으로 가장 흥미로운 순간에 집중적으로 빠진다. 통계에서는
이런 걸 \textbf{MNAR}(missing not at random, 무작위가 아닌 결측)이라고 부른다. 이 성질이
나중에 채점 방식 전체를 규정한다.

\subsection{질문}

\begin{note}
CES가 비어 있는 그 시점에서, \textbf{같은 순간의 빠른 진단(BES·ECEI·Mirnov)과 과거 CES
기록만으로}, CES 자기 시간축을 이어 붙이는 것만으로는 알 수 없는 \Ti\,·\,\Vrot\ 정보를
되찾을 수 있는가?
\end{note}

\subsection{답 (미리 말하면)}

\begin{itemize}
\item \textbf{\Ti: 된다.} 미래를 보는 보간을 이기고, 실시간으로 쓸 수 있는 가장 강한
  방법(인과 GP)도 이긴다. 4가지 분할 전부에서, 두 가지 데이터 처리(컷·포함) 전부에서.
\item \textbf{\Vrot: 안 된다.} 보간과 동률이다. 다만 \textbf{값이 급하게 변하는
  구간에서만은} 이긴다.
\item \textbf{왜 갈리는지도 밝혔다.} 모델 잘못이 아니라 \textbf{데이터에 회전을 보는 빠른
  신호가 없어서}다(\ref{sec:asym}장).
\end{itemize}

% =========================================================================
\section{데이터}

\subsection{무엇이 있나}

\begin{itemize}
\item \textbf{641개 방전} (shot 번호 30801--32751). 이건 KSTAR \textbf{2022년 실험 캠페인}이다.
\item 0.01초 격자, 총 \textbf{247{,}207행} (파일당 중앙값 339행)
\item 한 행에 들어 있는 것: BES 9개 + ECEI 4개 + Mirnov 2개 + 시각 + \Ti\ + \Vrot
\end{itemize}

한 방전 파일 안에서 측정이 0.5초 넘게 끊기면 \textbf{다른 구간(세그먼트)}으로 취급한다.
대부분의 파일은 \textbf{주 구간 하나}(중앙값 301행 $\approx$ 3.0초, 10--90퍼센타일
1.3--7.0초)에 $t=0$ 같은 자리에 떠 있는 고립된 행 몇 개가 붙은 모양이다.

\subsection{빈 곳이 얼마나 되나}

\begin{center}
\begin{tabular}{@{}lcc@{}}
\toprule
 & \Ti & \Vrot \\
\midrule
아예 빈 칸(NaN) & 8.2\,\% & 23.9\,\% \\
유지값(직전 값 복사) & 0.0\,\% & 41.1\,\% \\
\textbf{진짜 새 정보가 있는 칸} & \textbf{91.8\,\%} & \textbf{35.0\,\%} \\
값이 있는 칸 중 유지값 비율 & 0.0\,\% & \textbf{54.0\,\%} \\
\bottomrule
\end{tabular}
\end{center}

\subsection{데이터 품질 점검 \textcircled{1} --- 유지값}

\Vrot\ 값의 54\,\%가 유지값이고, 한 번 이어질 때 최대 \textbf{1{,}214행}(약 12초)까지 같은
값이 반복된다. 641개 파일 중 499개에 이 현상이 있다.

\textbf{혹시 진짜로 같은 값이 연속으로 측정된 건 아닐까?} 확인했다. \Vrot 는 소수점 5자리까지
기록되고 서로 다른 값들의 최소 간격이 0.00004라서, 우연히 완전히 같은 값이 연속으로 나올
가능성은 사실상 없다. \Ti 는 226{,}991행 중 딱 1행만 해당된다.

$\rightarrow$ \textbf{유지값은 학습·채점·정규화·보간 기준점 어디에서도 전부 제거}한다.

\subsection{데이터 품질 점검 \textcircled{2} --- 이온 온도 이상값}

관측된 \Ti 의 99퍼센타일은 2{,}089\,eV인데 최대값은 14{,}984\,eV다. 이 꼬리는 물리가 아니라
\textbf{분광 스펙트럼 맞춤 실패}다.

\begin{itemize}
\item 3{,}000\,eV 초과: \textbf{1{,}197행 (0.53\,\%)}, 951번의 연속 구간, 274개 방전에 걸침
\item 85\,\%는 딱 1행짜리고, 5행 이상 이어지는 경우는 2\,\%
\item 이상값의 크기는 이웃 평균의 \textbf{13배}(사분위 범위 6--26배)
\end{itemize}

\textbf{빼야 하나 넣어야 하나?} 둘 다 비판이 가능해서, \textbf{둘 다 1차 결과로 보고}하는
방식(\ref{sec:cut}절의 ``컷''/``포함'')을 택했다. 컷 문턱을 2.5 / 3 / 4\,keV로 바꿔 다시
학습해도 결과가 거의 같다는 것도 확인했다(\ref{sec:thresh}장).

한계도 적어둔다: 값으로 자르는 건 \textbf{한쪽 방향만} 걸러낸다. 위로 튀는 이상값의
19\,\%만 제거하고, \textbf{아래로 꺼지는 이상값 4{,}965행은 그대로 남는다}. 진짜 해법은
CES가 스펙트럼을 얼마나 잘 맞췄는지 알려주는 품질 지표를 받는 것이다(\ref{sec:data}장).

% =========================================================================
\section{우리가 만든 모델}

\subsection{왜 ``전체를 순서대로 읽는'' 방식인가}

이전 방식(윈도)은 \textbf{정답이 있는 행만 남기고} 그 앞 2줄을 잘라서 봤다. 그런데 생각해
보면 이상하다 --- \textbf{빠른 진단은 모든 칸에 다 있다}. 정답이 없다고 그 칸을 버릴 이유가
없다. 그래서 문제를 다시 세웠다:

\begin{note}
모든 칸을 시간 순서대로 읽되(정답이 없는 칸도 맥락으로 읽는다), 정답이 있는 칸에서만
오차를 계산한다.
\end{note}

이게 \textbf{``전체격자 시퀀스'' 방식}이고, 이 모델을 \textbf{\seqv}라고 부른다.

\subsection{\seqv 의 생김새 (파라미터 357{,}570개)}

한 칸마다 22개 숫자를 입력으로 받는다: 빠른 진단 15개 + 직전 칸과의 시간 간격 +
타겟별(마지막으로 잰 값, 그게 얼마나 오래됐는지, 잰 적이 있는지 여부).
그 위에 \textbf{서로 완전히 독립인 순환신경망(LSTM) 두 개}가 돈다:

\begin{itemize}
\item \textbf{\Ti\ 갈래} (2층, 폭 160): 위 22개를 \textbf{전부} 본다
\item \textbf{\Vrot\ 갈래} (1층, 폭 64): 빠른 진단이 아닌 \textbf{7개만} 본다
\end{itemize}

\textbf{왜 \Vrot\ 갈래에서 빠른 진단을 끊었나}: 실험으로 확인해 보니 빠른 진단에는 회전
정보가 \textbf{전혀 없기 때문}이다(\ref{sec:asym}장). 넣어봐야 잡음만 는다.

\textbf{중요한 구현 포인트}: 이 차단은 \textbf{출력단이 아니라 인코더 단계에서} 해야 한다.
순환 상태를 두 갈래가 공유하면 출력 배선을 어떻게 하든 빠른 진단 정보가 \Vrot\ 쪽으로
새어든다. 우리는 실제로 검증했다 --- \textbf{빠른 진단 15채널 값을 아무렇게나 흔들어도
\Vrot\ 출력이 비트 단위로 똑같다.}

\textbf{두 번째 중요한 점 --- 도달거리}: 이 모델의 과거 시야는 고정된 2줄이 아니라
\textbf{그 구간 전체}다. 나중에 보겠지만, 이 차이가 가장 강한 실시간 상대(인과 GP)를
이기는 결정적 이유다(\ref{sec:gate}절).

\subsection{비교용 대조군 --- 윈도 모델}

이전 방식도 그대로 살려서 \textbf{완전히 같은 데이터·같은 규칙으로} 학습시켰다. 정답 칸 앞
2줄을 보고, BES/ECEI/Mirnov 각각을 작은 합성곱 신경망으로 읽고, 과거 CES 기록은 GRU로 읽은
다음, \textbf{그 타겟이 실제로 측정된 행에만} 주의(attention)를 주도록 강제한다.
파라미터 201{,}258개. 같은 행에서 나란히 채점할 수 있으므로 \textbf{paired 비교}가 가능하다.

\subsection{모델을 키워봐야 소용없었다}

두 차례의 통제 실험에서 살아남은 개선은 \textbf{주의(attention) 방식 하나뿐}이었다. 모델을
키우거나, 지름길 연결을 넣거나, 추출기를 더 붙이는 건 전부 효과가 없었다. 확정 프로토콜에서도
같은 교훈이 두 번 더 반복됐다.

$\rightarrow$ \textbf{병목은 모델이 아니라 입력 정보에 있다.} 그래서 이 논문은 모델 설명이
짧고 채점 설명이 길다.

% =========================================================================
\section{어떻게 채점했나}

이 절이 이 연구의 실질적인 핵심이다. \textbf{``잘 나왔다''를 믿을 수 있게 만드는 규칙들}이다.

\subsection{미리 정해놓고 시작했다 (사전등록)}

시험 데이터를 한 번이라도 들여다보면 그 뒤의 모든 판단이 오염된다. 그래서 \textbf{시험
데이터를 채점하기 전에} 다음을 문서로 못 박았다:

\begin{center}
\begin{tabularx}{\textwidth}{@{}l L@{}}
\toprule
\textbf{규칙} & \textbf{내용} \\
\midrule
PR1 & 대표 비교 상대는 PCHIP으로 한다 \\
PR2 & 보간이 구간 경계에서 이웃을 못 찾으면 persistence로 대체하고,
  \textbf{그 비율을 반드시 보고}한다 \\
PR3 & 시험 세트는 최소 15방전 · \Ti\ 3{,}000행 이상이어야 한다 \\
PR4 & 유의 판정은 방전 단위 재추출 1만 번의 95\,\% 신뢰구간이 0을 넘지 않을 때만 \\
추가 & 유지값 제거 · 윈도 2 · 파일당 500샘플 상한 · 두 모집단(컷/포함) ·
  \textbf{모델 채택 규칙을 시험 채점 전에 커밋} \\
\bottomrule
\end{tabularx}
\end{center}

\textbf{PR2가 특히 중요하다}: \Vrot 는 보간이 이웃을 못 찾는 비율이 \textbf{40--44\,\%}나
된다. 즉 \Vrot 의 ``PCHIP 대비'' 점수는 실제로는 5번 중 2번이 ``persistence 대비''다.
이 사실을 숨기지 않고 보고한다.

\subsection{시험 세트는 얼려뒀다}

방전 파일 단위로 학습 / 검증 / 시험을 나누고, \textbf{시험 세트는 모델 고르는 과정에서 단
한 번도 보지 않았다}. 기준 분할(seed 42)의 시험 규모는 \Ti\ 32{,}589행 · 96방전,
\Vrot\ 10{,}463행 · 60방전이다.

분할 방식을 정하는 난수(seed)를 42 / 1 / 7 / 123 네 가지로 두었는데, \textbf{1·7·123은
모델을 고를 때 쓰지 않았다}. 그래서 이 셋에서도 결과가 재현되면 그건 ``선택 과정 바깥에서의
독립 재현''이다.

\subsection{채점 기준을 완전히 똑같이 맞췄다}

모든 방법이 \textbf{똑같은 (파일, 행) 집합}에서, \textbf{똑같은 마스크}로 채점된다. 이게
지켜지는지 파일 키를 비트 단위로 대조해서 확인한다. 보간은 자기가 맞혀야 할 값 자신은 절대
쓰지 못하고, 구간 경계를 넘어가지 못한다.

\subsection{왜 방전 단위로 신뢰구간을 계산하나}

같은 방전 안의 이웃 행들은 서로 강하게 닮았다. 행 단위로 재추출하면 \textbf{실제보다 표본이
훨씬 많은 척}하게 되어 신뢰구간이 거짓으로 좁아진다. 그래서 \textbf{방전 전체를 통째로
재추출}한다.

그 대가로 \textbf{유효 표본 수 = 방전 수}(\Ti\ 약 96, \Vrot\ 60--66)가 되고, \textbf{이게
이 연구 전체의 통계적 한계선}이다. \Vrot 가 전역에서 판정되지 않는 근본 이유도 여기 있다.

% =========================================================================
\section{결과 \textcircled{1} --- 평소 조건에서}

\subsection{오차 크기 순위 (분할 seed 42, 컷 데이터)}

\begin{center}
\begin{tabularx}{\textwidth}{@{}l L r r@{}}
\toprule
\textbf{방법} & \textbf{무엇을 보나} & \textbf{\Ti\ 오차 (eV)} & \textbf{\Vrot\ 오차 (km/s)} \\
\midrule
\textbf{\seqv\ (우리 모델)} & 빠른 진단 + 과거 CES, 구간 전체 & \textbf{157.8} & \textbf{23.6} \\
윈도 대조군 & 빠른 진단 + 과거 CES 2줄 & 169.2 & 26.1 \\
인과 GP & 과거 CES 이웃 16개 & 164.3 & 28.8 \\
persistence & 마지막 잰 값 & 197.2 & 33.4 \\
국소 AR & 과거 CES만 & 472.2 & 51.0 \\
GP (미래 사용) & 과거 + 미래 CES & 153.8 & 24.7 \\
선형 보간 (미래 사용) & 과거 + 미래 CES & 169.8 & 29.0 \\
PCHIP (미래 사용) & 과거 + 미래 CES & 173.6 & 30.2 \\
\bottomrule
\end{tabularx}
\end{center}

읽는 법: \textbf{작을수록 좋다.} 우리 모델(157.8)은 실시간으로 쓸 수 있는 모든 방법보다
낫고, \textbf{미래를 보는 보간(169.8, 173.6)보다도 낫다.} 유일하게 대등한 건 미래를 보는
GP(153.8)뿐이다.

\subsection{대표 결과 --- \Ti 는 두 데이터 처리 모두에서 4/4}

\begin{center}
\begin{tabular}{@{}llccc@{}}
\toprule
\textbf{타겟} & \textbf{분할} & \textbf{컷: PCHIP 대비 (95\,\% 구간)} &
  \textbf{포함: PCHIP 대비} & \textbf{인과 GP 대비} \\
\midrule
\Ti & 42 & \textbf{+0.174} (+0.097, +0.236) & \textbf{+0.225} (+0.109, +0.293) & \textbf{+0.078 / +0.154} \\
\Ti & 1 & \textbf{+0.248} (+0.188, +0.295) & \textbf{+0.238} (+0.153, +0.302) & \textbf{+0.133 / +0.169} \\
\Ti & 7 & \textbf{+0.257} (+0.199, +0.302) & \textbf{+0.292} (+0.232, +0.344) & \textbf{+0.138 / +0.123} \\
\Ti & 123 & \textbf{+0.264} (+0.188, +0.320) & \textbf{+0.316} (+0.186, +0.392) & \textbf{+0.105 / +0.149} \\
\midrule
\Vrot & 42 & \textbf{+0.390} (+0.077, +0.591) & \textbf{+0.384} (+0.066, +0.586) & \textbf{+0.331 / +0.324} \\
\Vrot & 1 & +0.183 ($-$0.028, +0.280) & \textbf{+0.195} (+0.000, +0.286) & +0.130 / \textbf{+0.143} \\
\Vrot & 7 & +0.135 ($-$0.358, +0.269) & +0.132 ($-$0.362, +0.266) & +0.020 / +0.018 \\
\Vrot & 123 & +0.305 ($-$0.049, +0.437) & +0.304 ($-$0.065, +0.433) & \textbf{+0.134} / +0.132 \\
\bottomrule
\end{tabular}
\end{center}

\textbf{표 읽는 법}: 괄호 안은 95\,\% 신뢰구간이다. \textbf{구간의 왼쪽 끝이 0보다 크면}
``우연이 아니다''로 판정하고 굵게 표시했다.

\begin{itemize}
\item \textbf{\Ti}: 8칸(4분할 $\times$ 2처리) \textbf{전부 성공}. 인과 GP도 8칸 전부
  이긴다(+0.08--+0.17). persistence는 +0.36--+0.46으로 크게 이긴다.
  $\rightarrow$ 논문의 \textbf{조건 없는 주장}이 된다.
\item \textbf{\Ti\ vs 미래를 보는 GP}: 동률이다($-$0.05--+0.11, 8칸 중 1칸만 유의).
  그래서 주장을 정확히 쓰면 이렇다: ``미리 정해둔 보간과 \textbf{실시간으로 쓸 수 있는 모든
  방법}을 이기고, \textbf{미래를 보며 표본마다 스스로 조율하는 최강 평활기와는 대등하다}.''
\item \textbf{\Vrot}: 점추정은 8칸 전부 양수지만, 신뢰구간 기준으로는 컷 1/4 · 포함 2/4다.
  \textbf{잡음과 구별되지 않으므로 ``회전도 이겼다''는 주장은 하지 않는다.}
\end{itemize}

\textbf{``포함'' 쪽 점수가 더 높은 이유}: 이상값을 남겨두면 \textbf{모든 방법이} PCHIP 대비
더 좋아 보인다(우리 모델 +0.268 vs +0.236, 대조군 +0.238 vs +0.173). 이상값이 보간의
기준점을 오염시키기 때문이다. 즉 ``포함'' 점수에는 \textbf{이상값에 얼마나 안 흔들리나}라는
별개 성분이 섞여 있다. 이게 한쪽 숫자만 인용하면 안 되는 이유다.

\subsection{전체를 읽는 방식이 정말 나은가 (모델 선정 관문)}
\label{sec:gate}

시험 데이터를 열기 전에 통과 조건 4개를 미리 정해놓고 확인했다. 분할 seed 4개 $\times$
초기값 seed 4개 = \textbf{16번 학습}했고, 각 결과를 자기 분할의 윈도 대조군과 \textbf{같은
행에서 나란히} 비교했다.

\begin{itemize}
\item \Ti: \textbf{16번 전부 양수, 그중 13번 유의}
\item 분할별 평균 +0.129 / +0.059 / +0.078 / +0.058 $\rightarrow$ \textbf{초기값에 따른
  흔들림보다 분할에 따른 흔들림이 훨씬 크다}
\item 전체 합산 \textbf{+0.081} (신뢰구간 +0.067--+0.096)
\item 학습 시간을 똑같이 맞춰도(고정 10에폭) 4/4 전부 양수 유지
\item \Vrot: 유의하게 나빠진 경우 \textbf{0번} (유의하게 좋아진 경우 8번)
\end{itemize}

\textbf{이 방식이 사주는 것이 정확히 무엇인가}: 윈도 대조군은 인과 GP와 대등할 뿐이다(1/4).
반면 시퀀스 모델은 인과 GP를 4/4+4/4로 이긴다. 차이는 \textbf{과거를 얼마나 멀리
보느냐(도달거리)}다. 그리고 이 방식은 \textbf{학습 비용이 1/10}이다 --- 윈도를 조립하고
부분집합을 늘리는 작업이 없기 때문이다.

\subsection{빈칸이 길수록 어떤가}

\begin{center}
\begin{tabular}{@{}llccc@{}}
\toprule
\textbf{빈칸 길이} & \textbf{행 수 (컷/포함)} & \textbf{방전} & \textbf{컷: PCHIP 대비} &
  \textbf{포함: PCHIP 대비} \\
\midrule
\Ti\ $\le$ 15\,ms & 134{,}546 / 135{,}317 & 301 & \textbf{+0.239} (+0.197, +0.274) & \textbf{+0.299} (+0.244, +0.347) \\
\Ti\ $>$ 15\,ms & 3{,}422 / 3{,}334 & 265 / 263 & \textbf{+0.268} (+0.187, +0.337) & \textbf{+0.206} (+0.108, +0.290) \\
\Ti\ $>$ 45\,ms & 460 / 429 & 104 / 101 & \textbf{+0.267} (+0.092, +0.414) & $-$0.004 ($-$0.304, +0.246) \\
\Vrot\ $\le$ 15\,ms & 51{,}689 & 197 & \textbf{+0.233} (+0.020, +0.318) & \textbf{+0.233} (+0.020, +0.317) \\
\Vrot\ $>$ 15\,ms & 466 / 456 & 130 & \textbf{+0.418} (+0.104, +0.680) & \textbf{+0.432} (+0.128, +0.696) \\
\bottomrule
\end{tabular}
\end{center}

\Ti 의 우세는 \textbf{바로 옆 칸에만 국한되지 않는다} --- 15\,ms 넘게 떨어진 곳에서도
유지된다.

\textbf{그리고 여기서 \Vrot 의 유일한 조건 없는 성공이 나온다}: 보간이 가장 힘들어하는
\textbf{15\,ms 넘는 빈칸}에서 \Vrot 도 PCHIP을 양쪽 데이터에서 이긴다(+0.418, +0.432).

% =========================================================================
\section{결과 \textcircled{2} --- 일부러 불리한 조건에서 다시 채점 (스트레스 테스트)}

\subsection{``스트레스 테스트''가 뭔가}

5장의 결과는 \textbf{평소 조건}에서 나온 것이다. 데이터를 무작위로 나누고, 값이 남아 있는
칸에서 채점했다. 그런데 이 조건은 \textbf{실제로 쓸 때의 조건과 다르다}. 두 가지가 다르다:

\begin{center}
\begin{tabularx}{\textwidth}{@{}L L L@{}}
\toprule
\textbf{평소 조건} & \textbf{실제 조건} & \textbf{그래서 뭘 확인해야 하나} \\
\midrule
\textbf{값이 있는 칸}을 일부러 가려서 채점 & 진짜로 \textbf{비어 있는 칸}을 채워야 함 &
  빈칸은 어려운 순간에 몰려 있는데(MNAR), 그래도 이기나? \\
방전을 \textbf{무작위로} 학습/시험에 나눔 & 과거 실험으로 배워서 \textbf{미래 실험}에 씀 &
  시간이 지나 장비 상태가 바뀌어도 이기나? \\
\bottomrule
\end{tabularx}
\end{center}

\textbf{스트레스 테스트 = 이 두 가지를 일부러 불리하게 바꿔놓고 다시 채점하는 것.} 평소
조건에서 잘 나온 결과가 여기서 무너지면, 그건 실제로는 못 쓰는 결과다.

이 절이 이 논문에서 가장 중요한 부분이다. \textbf{바로 여기서 이전 모델(윈도)이 무너졌고,
새 모델은 살아남았다.}

\subsection{스트레스 \textcircled{1} --- 진짜 빈칸의 분포에 맞춰 다시 채점 (MNAR 보정)}

\textbf{문제}: 진짜 빈칸에는 정답이 없으니 채점할 수가 없다. 그래서 \textbf{값이 있는 칸을
일부러 가려서} 채점한다. 그런데 빈칸은 어려운 순간에 몰려 있으므로, 이 방식은 \textbf{실제보다
후하게} 나온다.

\textbf{해법}: 값이 있는 칸과 빈칸 \textbf{양쪽에서 모두 계산할 수 있는 두 가지 특징}으로
데이터를 층으로 나눈다. (1) 마지막으로 잰 이후 얼마나 지났나 (15 / 25 / 45\,ms),
(2) 입력만 보고 판단한 활동 수준. 그러고 나서 \textbf{진짜 빈칸이 각 층에 얼마나 몰려
있는지의 비율로 점수를 다시 가중평균}한다. 이걸 \textbf{재가중}이라고 한다.

\textbf{얼마나 커버되나 (중요한 한계)}:
\begin{itemize}
\item 빈 \Ti\ 중 우리 모델을 적용할 수 있는(2줄 이내에 과거 관측이 있는) 비율:
  \textbf{54--68\,\%}
\item 빈 \Vrot 는 \textbf{4--6\,\%}뿐 --- 회전은 한 번 끊기면 아주 길게 끊기기 때문
\item $\rightarrow$ \textbf{재가중한 \Vrot\ 결과는 전체 빈칸의 1/20에 대한 답이라 결론으로
  쓰지 않는다.}
\end{itemize}

\begin{center}
\begin{tabular}{@{}llcccc@{}}
\toprule
\textbf{데이터} & \textbf{분할} & \multicolumn{2}{c}{\textbf{PCHIP 대비}} &
  \multicolumn{2}{c}{\textbf{persistence 대비}} \\
\cmidrule(lr){3-4}\cmidrule(lr){5-6}
 & & 보정 전 & \textbf{보정 후} & 보정 전 & \textbf{보정 후} \\
\midrule
컷 & 42 & +0.174 & +0.140 ($-$0.071, +0.290) & +0.360 & \textbf{+0.398} (+0.278, +0.518) \\
컷 & 1 & +0.248 & \textbf{+0.164} (+0.062, +0.250) & +0.406 & \textbf{+0.366} (+0.299, +0.448) \\
컷 & 7 & +0.257 & +0.203 ($-$0.024, +0.346) & +0.403 & \textbf{+0.310} (+0.227, +0.398) \\
컷 & 123 & +0.264 & \textbf{+0.283} (+0.069, +0.392) & +0.415 & \textbf{+0.381} (+0.246, +0.464) \\
\midrule
포함 & 42 & +0.225 & \textbf{+0.140} (+0.033, +0.250) & +0.423 & \textbf{+0.443} (+0.141, +0.623) \\
포함 & 1 & +0.238 & \textbf{+0.217} (+0.086, +0.319) & +0.436 & \textbf{+0.383} (+0.273, +0.536) \\
포함 & 7 & +0.292 & \textbf{+0.167} (+0.067, +0.265) & +0.406 & \textbf{+0.283} (+0.172, +0.380) \\
포함 & 123 & +0.316 & \textbf{+0.221} (+0.050, +0.320) & +0.459 & \textbf{+0.337} (+0.164, +0.455) \\
\bottomrule
\end{tabular}
\end{center}

\textbf{결과}:
\begin{itemize}
\item \textbf{실시간 시스템의 진짜 경쟁자인 persistence 대비로는 8칸 전부 살아남는다}
  (+0.28--+0.44). 보정 때문에 깎이는 폭은 최대 0.12에 불과하다.
\item 미래를 보는 PCHIP 대비로는 점추정(+0.14--+0.28)은 유지되지만, 신뢰구간이 컷 데이터
  2개 분할에서 0을 넘는다(컷 2/4, 포함 4/4).
\end{itemize}

\textbf{정확한 문장}: ``진짜로 비어 있고 모델을 적용할 수 있는 시점에서, 우리 모델은
\textbf{실시간으로 쓸 수 있는 모든 방법보다 확실히 낫고}, 미래를 보는 보간보다는
\textbf{데이터 처리 방식에 따라 조건부로} 낫다.''

\subsection{스트레스 \textcircled{2} --- 과거 실험으로 배워 미래 실험에 쓰기 (캠페인 분할)}

\textbf{문제}: 무작위로 나누면 학습에 쓴 방전과 시험에 쓴 방전이 \textbf{같은 날 것}일 수
있다. 실제 배치 상황은 다르다 --- 지금까지 쌓인 데이터로 배워서 \textbf{앞으로 할 실험}에
쓴다. 그사이 장비 교정도 바뀌고 벽 상태도 바뀐다.

\textbf{해법}: 641개 방전을 shot 번호 순으로 줄 세워 \textbf{시간으로만} 자른다.
학습 416개 (30801--31991) / 검증 128개 (32002--32310) / 시험 97개 (32312--32751).
즉 \textbf{어떤 시험 방전도 학습 방전보다 앞서지 않는다.}

\begin{center}
\footnotesize
\setlength{\tabcolsep}{4.5pt}
\begin{tabular}{@{}llccc@{}}
\toprule
\textbf{데이터} & \textbf{방법} & \textbf{\Ti\ PCHIP 대비 (초기값 4개)} & \textbf{성공} &
  \textbf{인과 GP} \\
\midrule
컷 & 윈도 대조군 & +0.027 / \textbf{+0.091} / $-$0.001 / \textbf{+0.061} & 2/4 & \textbf{0/4} \\
컷 & 윈도 + shot별 표준화 & \textbf{+0.103 / +0.107 / +0.094 / +0.107} & 4/4 & --- \\
컷 & \textbf{\seqv} & \textbf{+0.187 / +0.174 / +0.181 / +0.177} & \textbf{4/4} & \textbf{4/4} \\
\midrule
포함 & 윈도 대조군 & +0.014 / +0.047 / +0.055 / +0.089 & \textbf{0/4} & \textbf{0/4} \\
포함 & \textbf{\seqv} & \textbf{+0.173 / +0.202 / +0.198 / +0.184} & \textbf{4/4} & \textbf{4/4} \\
\bottomrule
\end{tabular}
\end{center}

\seqv 가 같은 분할의 윈도 대조군을 같은 행에서 얼마나 앞서는지(paired):
컷 \textbf{+0.164 / +0.091 / +0.182 / +0.124}, 포함
\textbf{+0.161 / +0.163 / +0.151 / +0.104} --- \textbf{8칸 전부}.

\textbf{여기가 이 논문의 결정적 장면이다.}

\begin{itemize}
\item \textbf{윈도 모델은 무너진다.} 미래를 보는 보간 대비 우위가 사라지고(2/4 · 0/4),
  인과 GP는 아예 못 이긴다(0/4).
\item \textbf{시퀀스 모델은 안 무너진다.} PCHIP 4/4+4/4, 인과 GP 4/4+4/4, 윈도 대조군은
  8/8로 이긴다.
\end{itemize}

\textbf{왜 무너졌는지도 측정했다} (짐작이 아니다):

\begin{center}
\begin{tabular}{@{}lc@{}}
\toprule
 & \textbf{학습 구간 $\rightarrow$ 시험 구간 이동량} \\
\midrule
BES (밀도 요동 센서) & \textbf{1.22$\sigma$} (크기 비 0.75배) \\
ECEI (전자 온도 센서) & 0.53$\sigma$ (크기 비 0.62배) \\
타겟(\Ti\,·\,\Vrot) & \textbf{0.115$\sigma$} (크기 비 1.06배) \\
\bottomrule
\end{tabular}
\end{center}

즉 \textbf{물리가 변한 게 아니라 센서 눈금이 변했다.} 타겟은 거의 그대로인데 입력만
5--11배 크게 흘렀다. 학습 데이터 전체로 정규화 기준을 잡으면, 시간이 지나 눈금이 바뀐
데이터에서 그 기준이 어긋난다.

\textbf{고치는 법도 확인했다}: 빠른 진단을 \textbf{각 방전 내부에서 표준화}하면(타겟은
그대로) 윈도 모델이 2/4 $\rightarrow$ 4/4로 회복한다. 그리고 \Vrot 는 전혀 변하지
않는다($-$0.003--+0.008) --- \textbf{\Vrot\ 갈래는 빠른 진단을 아예 안 보므로 변하면 안
되고, 실제로 안 변했다.} 이건 우리 구조가 의도대로 작동한다는 \textbf{음성 대조 실험}이다.

\seqv 는 이 표준화를 정의상 이미 하고 있고, 거기에 도달거리까지 갖고 있다. 그래서 애초에
안 무너진다.

\subsection{두 스트레스 요약}

\begin{center}
\begin{tabularx}{\textwidth}{@{}L L L@{}}
\toprule
\textbf{채점 조건} & \textbf{미래 보는 보간(PCHIP) 대비} & \textbf{실시간 방법 대비} \\
\midrule
평소 (무작위 분할, 값 있는 칸) & \textbf{4/4 · 4/4}, +0.17--+0.32 &
  \textbf{4/4 · 4/4} vs 인과 GP, +0.08--+0.17 \\
진짜 빈칸 분포로 보정 & 2/4 · \textbf{4/4}, +0.14--+0.28 &
  \textbf{4/4 · 4/4} vs persistence, +0.28--+0.44 \\
미래 실험으로 시간 분할 & \textbf{4/4 · 4/4}, +0.17--+0.20 &
  \textbf{4/4 · 4/4} vs 인과 GP, +0.11--+0.16 \\
\bottomrule
\end{tabularx}
\end{center}

남는 주의 두 가지: 캠페인 실험의 ``4번''은 \textbf{분할 4가지가 아니라 한 시간 분할 위의
초기값 4가지}다. 그리고 컷 데이터에서 4번 중 2번이 최대 에폭(30)까지 가서 멈췄다.

% =========================================================================
\section{왜 이온 온도는 되고 회전은 안 되나}
\label{sec:asym}

\textbf{방법}: 입력의 일부를 일부러 지우고 다시 채점한다(절제 실험). 무엇을 지웠을 때
성능이 무너지는지가 곧 ``그 정보가 어디서 왔는지''를 말해준다.

\begin{center}
\begin{tabular}{@{}llll@{}}
\toprule
\textbf{지운 것} & \textbf{타겟} & \textbf{컷} & \textbf{포함} \\
\midrule
아무것도 안 지움 & \Ti & +0.173 & +0.238 \\
\textbf{빠른 진단 지움} & \Ti & \textbf{$-$0.125} (변화 $-$0.25--$-$0.43, 4/4 유의) & +0.201 (변화 $-$0.03--$-$0.09) \\
\textbf{과거 CES 지움} & \Ti & \textbf{$-$2.11} (변화 $-$1.8--$-$4.6, 4/4 유의) & $-$1.16 (4/4 유의) \\
\midrule
아무것도 안 지움 & \Vrot & +0.213 & +0.206 \\
\textbf{빠른 진단 지움} & \Vrot & +0.213 · \textbf{변화 정확히 0.000} & +0.206 · \textbf{변화 정확히 0.000} \\
\textbf{과거 CES 지움} & \Vrot & \textbf{$-$2.89} (변화 $-$1.8--$-$6.3, 4/4 유의) & $-$3.51 (4/4 유의) \\
\bottomrule
\end{tabular}
\end{center}

\textbf{읽어야 할 것 세 가지}:
\begin{enumerate}
\item \textbf{과거 CES 기록은 두 타겟 모두에 필수다.} 지우면 성능이 $-$1에서 $-$4까지
  곤두박질친다.
\item \textbf{\Ti 가 보간을 이기는 그 마진은 빠른 진단에서 온다.} 빠른 진단을 지우면
  보간보다 못해진다($-$0.125). 물리적으로 당연하다 --- 전자와 이온은 충돌로 열을 주고받으므로,
  ECEI가 보는 전자 온도와 BES가 보는 밀도가 이온 온도를 제약한다.
\item \textbf{\Vrot 는 빠른 진단에서 정보를 전혀 못 얻는다.} 빠른 진단 값을 아무리 흔들어도
  출력이 \textbf{비트 하나까지 똑같다}. 이건 구조적으로 그렇게 만든 것이기도 하지만, 애초에
  그렇게 만든 이유가 실험 결과였다.
\end{enumerate}

\textbf{왜 회전 정보가 없나 --- 두 가지 이유}:
\begin{itemize}
\item \textbf{NBI 토크 채널이 데이터에 없다.} 플라즈마를 돌리는 힘은 중성입자빔이 주는
  회전력인데, 그 값이 데이터셋에 아예 없다.
\item \textbf{Mirnov가 망가져 있다.} 원래 kHz로 진동하는 자기 신호를 필터 없이 100\,Hz로
  줄여 저장하는 바람에 정보가 파괴됐다(\ref{sec:data}장에서 자세히).
\end{itemize}

\textbf{보너스로 알게 된 것}: ``포함'' 데이터에서는 \textbf{과거 CES만 쓰는 모델도 PCHIP을
+0.15--+0.23으로 이긴다.} 빠른 진단이 더해주는 건 0.03--0.09뿐이다. 즉 ``포함'' 점수에는
\textbf{이상값에 안 흔들리는 능력}이 섞여 있다. 빠른 진단의 순수 기여를 보려면 ``컷''
데이터를 봐야 한다. \textbf{두 데이터를 함께 보고해야 하는 이유가 이렇게 측정으로 나온다.}

% =========================================================================
\section{과거 기록은 얼마나 필요한가 --- 딱 하나면 된다}

과거를 몇 줄까지 보게 할지($W$)를 정하기 위해, $W$를 2·3·4·6·8로 바꿔가며 각각 4번씩 =
\textbf{24번 학습}했다. 여기에 \textbf{과거 CES를 아예 0으로 만든 조건}도 4번 추가했다.

\begin{center}
\begin{tabular}{@{}llcccc@{}}
\toprule
\textbf{쓸 수 있는 과거 관측} & \textbf{$W$} & \textbf{\Ti\ 점수} & \textbf{성공} &
  \textbf{\Vrot\ 점수} & \textbf{성공} \\
\midrule
\textbf{0개} & 4 (과거 차단) & \textbf{$-$0.026} & 0/4 & \textbf{$-$0.783} & 0/4 \\
\textbf{1개} & 2 & +0.238 & \textbf{4/4} & \textbf{+0.206} & 0/4 \\
2개 & 3 & \textbf{+0.246} & \textbf{4/4} & +0.203 & 1/4 \\
3개 & 4 & +0.221 & 3/4 & +0.190 & 1/4 \\
5개 & 6 & +0.190 & 3/4 & +0.205 & 1/4 \\
7개 & 8 & +0.216 & \textbf{4/4} & +0.204 & 2/4 \\
\bottomrule
\end{tabular}
\end{center}

\textbf{두 가지가 보인다}:
\begin{enumerate}
\item \textbf{과거 기록은 필수다.} 없애면 \Ti 는 보간보다 못해지고($-$0.026), \Vrot 는
  보간 오차의 1.8배로 붕괴한다($-$0.783).
\item \textbf{하나만 있으면 끝이다.} 과거 관측 1개를 주는 순간 두 타겟이 동시에 최고 수준으로
  올라오고, 그다음부터는 \textbf{평평하다}. \Ti 는 0.190--0.246, \Vrot 는 0.190--0.206
  범위인데, \textbf{같은 조건 안에서 난수만 바꿔도 0.07--0.16씩 흔들린다.} 즉 곡선 전체의
  변화보다 잡음이 더 크다.
\end{enumerate}

$\rightarrow$ ``평평해지는 지점 중 가장 작은 값''이라는 규칙에 따라 \textbf{$W = 2$}를 택했다.

\textbf{넓은 윈도를 쓸 유일한 이유는 정확도가 아니라 적용 범위였다.} 윈도가 좁으면 근처에
관측이 없는 칸은 아예 채점 대상에서 빠진다(15\,ms 넘는 빈칸 채점 표본이 1{,}958
$\rightarrow$ 456으로 줄어든다). 그런데 이건 \textbf{긴 윈도의 논거가 아니라 시퀀스 방식의
논거다} --- 시퀀스 모델은 도달거리가 구간 전체라서 $W$라는 설정값 자체가 없다.

% =========================================================================
\section{얼마나 단순하게 만들 수 있나}

\textbf{질문}: 35만 개 숫자짜리 모델이 하는 일을, 훨씬 작고 \textbf{사람이 읽을 수 있는}
모델로 재현할 수 있나?

\textbf{만든 것 (\bthree, 21{,}498개 숫자)}: \seqv 의 구조와 배선은 그대로 두고, 마지막
출력부만 이렇게 바꿨다.

\begin{note}
\textbf{예측 = 마지막으로 잰 값 + (작은 숫자 8개 각각에 가중치를 곱해 더한 보정) + 상수}
\end{note}

여기서 ``작은 숫자 8개''는 $-1$에서 $1$ 사이로 눌린 값이고, 가중치는 \textbf{0에서
시작}한다. 즉 \textbf{학습 시작 시점의 이 모델은 정확히 ``직전 값 그대로''(persistence)와
같다.} 학습이 진행되면서 보정을 배운다. 어떤 방전이든 ``직전 값이 이만큼이었고, 8개 항이
각각 이만큼 보정했다''로 완전히 분해해서 설명할 수 있다.

\begin{center}
\begin{tabular}{@{}lrcc@{}}
\toprule
\textbf{방법} & \textbf{숫자 개수} & \textbf{\Ti\ 점수 (컷)} & \textbf{\Ti\ 점수 (포함)} \\
\midrule
persistence & 0 & $-$0.264 & $-$0.288 \\
이름 붙은 항 3개 (anchor+$\Delta$) & 1{,}258 & $-$0.261 & $-$0.287 \\
\textbf{\bthree} & \textbf{21{,}498} & \textbf{+0.237} & +0.126 \\
윈도 대조군 & 201{,}258 & +0.173 & +0.238 \\
\textbf{\seqv} & \textbf{357{,}570} & \textbf{+0.236} & \textbf{+0.268} \\
\bottomrule
\end{tabular}
\end{center}

\textbf{결과 \textcircled{1}}: \textbf{``컷'' 데이터에서는 21{,}498개짜리가 357{,}570개짜리와
완전히 대등하다} (차이 평균 +0.002, 신뢰구간이 전부 0을 포함). 17배 작은 모델이 같은 일을 한다.

그 8개 숫자가 뭘 담고 있는지도 확인했다(선형 탐침 $R^2$): 직전에 잰 \Ti\ 값
\textbf{0.47--0.75}, ECEI가 보는 전자 온도 \textbf{0.31--0.48}, 국소 활동 수준 0.09--0.13
(거의 안 담음), 얼마나 오래됐는지 0.01--0.07 (거의 안 담음). 보정 항이 예측 분산의
25--39\,\%를 설명한다.

\textbf{결과 \textcircled{2}}: \textbf{``포함'' 데이터에서는 대등하지 않다}
($-$0.16--$-$0.21, 4/4 유의). 이유가 명확하다 --- 보정 폭이 $-1$--$1$로 \textbf{묶여
있어서}, 직전 값이 이상값(수천 eV)일 때 그걸 되돌릴 수가 없다. 실제로 ``직전 값 오차가
2\,keV 넘는 행''은 시험 데이터의 0.6--1.3\,\%에 불과한데 \textbf{이 모델 전체 오차의
73--83\,\%}를 차지한다(다른 모든 방법도 70--83\,\%로 비슷하다).

\textbf{결과 \textcircled{3} --- 모델 크기를 키우는 축은 닫혔다}: \Ti\ 인코더 폭을 24에서
260까지(파라미터 34k / 49k / 114k / 358k / 879k) 바꿔가며 재봤다. 점수는 +0.230 / +0.236 /
+0.235 / +0.236 / +0.230. \textbf{26배를 키워도 $\pm$0.008 안이다.}

$\rightarrow$ \textbf{성능의 한계는 모델이 아니라 정보에 있다.} 남은 개선은 전부 데이터
쪽이다(\ref{sec:data}장).

% =========================================================================
\section{어디서 이기고 있나 --- 급변 구간}

\textbf{``peak''와 ``본류''}: 값이 급하게 변하는 구간을 \textbf{peak}, 완만하게 흐르는
구간을 \textbf{본류}라고 부른다. peak 판정은 \textbf{답으로 쓸 값 자신을 빼고} 이웃 값들의
기울기와 분산만으로 하므로, ``정답을 보고 정답을 정의하는'' 순환 논리가 아니다.

\begin{center}
\begin{tabularx}{\textwidth}{@{}l L L@{}}
\toprule
\textbf{구간} & \textbf{\Ti\ (컷 / 포함)} & \textbf{\Vrot\ (컷 / 포함)} \\
\midrule
\textbf{peak (급변)} & \textbf{+0.45--+0.61 / +0.62--+0.72, 8/8 성공} &
  +0.54--+0.79 (8/8 양수, 각 2/4 유의) \\
본류 (완만) & +0.09--+0.20 (컷 4/4) / +0.06--+0.19 (포함 2/4) &
  거의 0 ($-$0.07--+0.15, 0/8) \\
\bottomrule
\end{tabularx}
\end{center}

\textbf{해석}: \textbf{보간은 매끄러운 구간에서는 이미 거의 최적이다.} 앞뒤 값을 이으면 그게
정답에 가깝다. 우리 모델의 가치는 \textbf{값이 급하게 움직이는 순간} --- 즉 물리적으로
중요한 순간에 있다.

\Vrot 도 마찬가지다. \textbf{전체로는 동률이지만 급변 구간에서만은 이긴다}(persistence
대비로는 +0.75--+0.86으로 8/8 성공). 비대칭은 전역이 아니라 \textbf{국소적}이다.

% =========================================================================
\section{이상값 자르는 기준을 바꿔보면}
\label{sec:thresh}

3{,}000\,eV라는 기준이 자의적이지 않은지 확인했다. 2{,}500 / 3{,}000 / 4{,}000\,eV로 각각
다시 학습했다.

\begin{itemize}
\item \Ti: +0.230 / +0.236 / +0.232 --- 사실상 같음. 유의 판정도 전부 4/4
\item \Vrot: +0.252 / +0.253 / +0.257 --- 사실상 같음
\end{itemize}

$\rightarrow$ \textbf{문턱값은 결과에 영향이 없다.} 중요한 건 ``자르느냐 마느냐''(두
모집단)이지 ``어디서 자르느냐''가 아니다.

% =========================================================================
\section{모델을 어떻게 골랐나}

시험 데이터를 보고 나서 모델을 고르면 그건 부정행위다. 그래서 규칙을 먼저 정했다.

\begin{itemize}
\item \textbf{윈도 계열}: 한 번에 한 가지만 바꾸고, 검증 데이터 점수로만 채택/기각
\item \textbf{주 모델 선정}: 통과 조건 4개를 미리 못 박고 확인 (\ref{sec:gate}절)
\item \textbf{그 뒤 후보 하나}(\seqv 에 주의 기법을 얹은 것): 4/4 전부 양수였지만
  (+0.009--+0.033) 유의는 1/4뿐이라 \textbf{미리 정한 기준(3/4 이상)에 미달} $\rightarrow$
  \textbf{채택하지 않았다}
\end{itemize}

마지막 항목이 중요하다. \textbf{이 후보는 검증 데이터에서는 2/2 유의였다.} 검증 데이터를
기준으로 골랐다면 채택했을 것이고, 시험 데이터에서 재현되지 않았을 것이다. \textbf{채택
기준을 시험 데이터에 두는 이유가 바로 이것이다.}

% =========================================================================
\section{실제로 쓸 수 있나}

\subsection{얼마나 빠른가}

실시간 제어에 쓰려면 0.01초 안에 답이 나와야 한다. 측정 조건은 \textbf{상태를 이어가며
한 칸씩 처리}(은닉 상태를 다음 칸으로 넘김), 일반 노트북급 기계, 워밍업 후 1{,}000회 반복이다.

\begin{center}
\begin{tabular}{@{}lcc@{}}
\toprule
\textbf{방법} & \textbf{CPU 중앙값} & \textbf{CPU 최악(99\,\%)} \\
\midrule
\textbf{\seqv\ 한 칸} & \textbf{1.05\,ms} & \textbf{1.61\,ms} \\
윈도 대조군 ($W=2$) & 3.8\,ms & 18.9\,ms \\
윈도 대조군 ($W=4$) & 4.0\,ms & 8.1\,ms \\
\bottomrule
\end{tabular}
\end{center}

\textbf{10\,ms 예산의 16\,\%만 쓴다.} 즉 80\,\% 이상이 남는다.

GPU는 1.21 / 2.31\,ms로 \textbf{오히려 느리다}. 이 크기(35만 파라미터)에서 한 칸씩 처리하면
GPU 커널을 띄우는 오버헤드가 계산량보다 크기 때문이다. $\rightarrow$ \textbf{제어 컴퓨터의
CPU에서 상태를 이어가며 돌려라.}

\subsection{예측에 오차범위를 붙일 수 있나}

\textbf{split conformal}이라는 방법을 썼다. 원리는 단순하다 --- 검증 데이터에서 ``이 모델이
보통 이만큼 틀린다''를 재고, 그 값을 시험 데이터의 예측에 $\pm$로 붙인다. \textbf{모델을
다시 학습할 필요가 없다.} 목표 신뢰도는 90\,\%로 잡았다.

공정하게 비교하려고 \textbf{PCHIP과 persistence에도 똑같은 절차를 적용}했다. 그래야 차이가
``기법의 차이''가 아니라 ``구간 품질의 차이''가 된다.

\textbf{결과: 4분할 $\times$ 2데이터 $\times$ 2타겟 $\times$ 2방식 = 32칸 전부에서 우리
모델의 구간이 두 기준선보다 낫다.}

\begin{center}
\begin{tabular}{@{}lccc@{}}
\toprule
 & \textbf{우리 모델} & \textbf{PCHIP} & \textbf{persistence} \\
\midrule
\Ti\ 점수 (컷, 낮을수록 좋음) & \textbf{1{,}272} & 1{,}554 & 1{,}727 \\
\Ti\ 점수 (포함) & \textbf{2{,}290} & 2{,}851 & 3{,}120 \\
\Vrot\ 점수 & \textbf{150} & 164 & 179 \\
\bottomrule
\end{tabular}
\end{center}

재미있는 점: ``포함'' 데이터에서 우리 모델의 구간은 PCHIP보다 \textbf{더 넓은데도}(반폭
224--255 vs 211--241\,eV) 점수는 더 좋다. 이상값이 있을 때 \textbf{틀렸을 때의 벌점이
커지는데, 우리 모델이 덜 놓치기 때문}이다.

\textbf{한계}: 보장되는 건 \textbf{전체 평균 신뢰도}(\Ti\ 0.87--0.92, \Vrot\ 0.91--0.94)
뿐이다. \textbf{방전별로 보면 보장되지 않는다} --- 어떤 방전에서는 구간이 체계적으로 좁을
수 있다.

% =========================================================================
\section{남은 개선은 전부 데이터 쪽이다}
\label{sec:data}

9장에서 모델 크기 축이 닫혔으므로, 남은 이득은 데이터에서만 나온다. 세 가지다.

\subsection{CES 스펙트럼 맞춤 품질 지표}

지금은 ``3{,}000\,eV 넘으면 실패로 본다''는 \textbf{값 기준}으로 자른다. 이건 위로 튀는 것만
걸러내는 \textbf{한쪽짜리 대리 지표}다. CES가 스펙트럼을 얼마나 잘 맞췄는지(맞춤 오차,
신호 세기)를 직접 받으면 \textbf{두 모집단으로 나눌 필요 자체가 없어진다.} \Vrot\ 이상값도
같은 규칙으로 처리된다.

\subsection{Mirnov 원본 신호 --- 정보는 전처리가 파괴했다}

\textbf{증거}: 같은 구간 안에서 바로 다음 칸과의 상관을 재보면
BES \textbf{+0.568}(연속적), ECEI \textbf{+0.572}(연속적),
\textbf{Mirnov $-$0.009}(완전한 백색잡음, 구간의 82\,\%가 $|r| < 0.1$).

이건 \textbf{kHz로 진동하는 자기 신호를 반앨리어싱 필터 없이 100\,Hz로 솎아낸} 전형적인
흔적이다. 연속한 샘플끼리의 위상 관계가 무작위가 되어, \textbf{진폭과 모드 번호 정보가
모델에 닿기 전에 파괴됐다.} 솎아낸 신호를 아무리 가공해도(적분, 절대값, 이동 RMS 등)
복구되지 않는다는 것도 실험으로 확인했다.

\textbf{해야 할 일}: 원본 kHz 신호에서 \textbf{구간 RMS · 대역 파워 · 모드 번호 · 모드 회전
주파수}를 뽑아 \Vrot\ 갈래에 넣는다.

\textbf{모드 회전 주파수가 특히 중요하다.} MHD 모드가 도는 속도는 플라즈마 회전과 물리적으로
직결되는 \textbf{유일한 빠른 신호 후보}다. \Vrot 가 안 되는 이유가 ``회전을 보는 빠른 신호가
없어서''였으므로, 이게 채워지면 결론 자체가 바뀔 수 있다.

\subsection{NBI 토크 채널}

\textbf{증거}: 각 타겟을 20번 이상 잰 538개 방전에서 전자 온도 $\leftrightarrow$ 이온 온도는
\textbf{$r = +0.353$} ($p = 3 \times 10^{-17}$, 확실히 상관 있음)인 반면, 전자 온도
$\leftrightarrow$ 회전은 \textbf{$r = +0.024$} ($p = 0.58$, 상관 없음)이다.

인과 사슬의 앞부분(빔 $\rightarrow$ 전자 가열 $\rightarrow$ 전자 온도)은 데이터에서
확인되므로, 뒷부분이 무상관이라는 건 \textbf{정보다}. \textbf{파워는 토크가 아니다.} 토크는
빔 에너지·접선 반경·주입 각도에 따라 달라지고, 회전은 국소 전자 온도가 아니라 운동량 수송과
가장자리 제동이 지배한다.

% =========================================================================
\section{한계 (숨기지 않고 적는다)}

\begin{itemize}
\item \textbf{통계적 검정력}: 유효 표본이 방전 수(\Ti\ 96, \Vrot\ 60--66)로 제한된다.
  \Vrot 가 전역에서 판정되지 않는 근본 원인이다. ``포함'' 데이터에서는 약 1\,\%의 행이 전체
  오차의 70--83\,\%를 차지해 소수 방전이 결과를 흔든다.
\item \textbf{MNAR 보정의 도달 범위}: \Ti\ 54--68\,\%, \Vrot 는 4--6\,\%만 커버한다.
\item \textbf{오프라인 주장의 상한}: 미래를 보는 GP와는 대등할 뿐이다(8칸 중 1칸만 유의).
\item \textbf{값 기준 컷은 대리 지표}다(한쪽 방향만 걸러냄).
\item \textbf{캠페인 실험은 한 시간 분할 위의 초기값 4개}이지 분할 4개가 아니다. 컷
  데이터에서 4번 중 2번이 최대 에폭까지 갔다.
\item \textbf{shot별 표준화는 방전 전체 통계를 쓰므로 오프라인 형태}다. 실시간으로 쓸 수 있는
  형태(누적 윈도·지수이동평균)와의 격차는 아직 안 재봤다.
\item \textbf{conformal 신뢰도는 전체 평균}이지 방전별로는 보장되지 않는다.
\item \textbf{지연 시간은 신경망 계산만} 잰 것이다. 전원 상태에 따라 절대값이 2배까지 차이
  난다(순서는 안 바뀐다).
\item \textbf{단일 장치, 두 종류의 순환 신경망 계열}만 시험했다.
\end{itemize}

% =========================================================================
\section{결론}

\textbf{한 문장}: 빠른 진단과 성긴 과거 기록만으로 CES 빈칸을 채우는 모델을 만들었고,
\textbf{이온 온도는 미래를 보는 보간까지 이겼지만 회전은 이기지 못했으며, 왜 그런지도
밝혔다.}

\textbf{구체적으로}:
\begin{enumerate}
\item \Ti 는 \textbf{4가지 분할 $\times$ 2가지 데이터 처리 전부}에서 미리 정한 보간을
  유의하게 이긴다(+0.17--+0.32). 실시간으로 쓸 수 있는 가장 강한 방법(인과 GP)은 8칸 전부에서
  이긴다. 미래를 보는 최강 평활기(GP)와는 대등하다.
\item \textbf{두 스트레스를 다 견딘다.} 진짜 빈칸 분포로 보정해도 실시간 방법 대비 8/8 생존,
  미래 실험으로 시간 분할해도 4/4+4/4. \textbf{대체된 윈도 모델은 시간 분할에서 완전히
  무너졌고}, 그 원인(빠른 진단이 5--11배 흐름)과 처방(방전별 표준화)까지 측정으로 특정했다.
\item \textbf{안 되는 것도 그대로 보고한다.} 회전은 전역 동률이고(급변 구간에서만 이김),
  값 기준 컷은 한쪽짜리 대리 지표이며, 빈칸 보정은 \Ti 의 절반 · \Vrot 의 1/20만 닿고,
  오차범위는 방전별로는 보장되지 않는다.
\item \textbf{모델을 키우는 방향은 닫혔다.} 21{,}498개짜리 해석 가능한 모델이 ``컷''
  조건에서 35만 개짜리와 대등하고, 폭을 26배 키워도 평평하다. \textbf{남은 이득은 전부
  데이터에 있다} --- CES 품질 지표, 원본 kHz Mirnov, NBI 토크 채널.
\end{enumerate}

% =========================================================================
\appendix
\section{표기와 재현}

\begin{itemize}
\item \texttt{skill} $= 1 - ($우리 오차$^2) / ($상대 오차$^2)$. 양수 = 우리가 나음.
\item ``유의'' = 방전 단위 재추출(1만 번) 95\,\% 신뢰구간이 0을 넘지 않음.
  ``n/4'' = 4개 분할 중 유의한 개수.
\item ``컷 / 포함'' = \Ti\ 이상값을 뺀 데이터 / 넣은 데이터. \textbf{둘 다 1차 결과.}
\item 코드: 주 모델 \texttt{experiments/seq/model\_seq\_v2.py}, 해석 가능 모델
  \texttt{model\_seq\_b3.py}, 윈도 대조군 \texttt{model\_iter009.py}(해시로 고정).
\item 수치 흐름: \texttt{collect\_paper\_numbers.py} $\rightarrow$
  \texttt{paper\_numbers.json} $\rightarrow$ \texttt{main.tex} / \texttt{main\_ko.tex} /
  \texttt{make\_figures\_en.py}.
\end{itemize}

\section{남은 일}

\begin{itemize}
\item[$\square$] 발표 자료 5종을 새 수치로 다시 만들기 (현재 옛 수치 상태)
\item[$\square$] 원본 kHz Mirnov 특징이 도착하면 규칙을 먼저 커밋한 뒤 소규모 시험
  $\rightarrow$ 확대
\item[$\square$] 승상님 검토: 초록·기여 문구, 두 모집단 서술 톤, 제목 유지 여부
\end{itemize}

\textbf{결정 기록 (2026-08-16)}:
\textcircled{1}~두 모집단(컷/포함) 공동 1차 유지 \quad
\textcircled{2}~\Vrot\ 처리 규칙 불변 \quad
\textcircled{3}~kHz Mirnov 미도착으로 대기

\end{document}
